\documentclass[12pt, a4paper]{article}

% --- Pacchetti Scientifici e Formattazione ---
\usepackage{amsmath, amssymb, amsfonts, amsthm}
\usepackage{booktabs, graphicx}
\usepackage[document]{ragged2e}
\usepackage{xcolor}

% --- Codifica e Font ---
\usepackage[utf8]{inputenc}
\usepackage[T1]{fontenc}
\usepackage[italian]{babel}
\usepackage{mathptmx}

% --- Gestione Margini e Layout ---
\usepackage[includeheadfoot]{geometry}
\geometry{
    a4paper,
    top=3cm,
    bottom=3cm,
    left=3.5cm,
    right=3.5cm,
}

% --- Gestione Interlinea e Figure ---
\usepackage{setspace}
\setstretch{1.3}
\usepackage[figuresright]{rotating}
\usepackage[pdftex, hidelinks]{hyperref}

\usepackage{subfiles}

% Document metadata
\title{Gestione Interviste}
\author{Mattia Cavina}
\date{09 febbraio 2026}

\begin{document}

\subsection{Raccolta dei requisiti}

\subsubsection{Periodo}
Le interviste si sono svolte dal 28 dicembre 2025 al 5 febbraio 2026.

\subsubsection{Stakeholder intervistati}
La raccolta dei requisiti è stata condotta mediante interviste con utenti finali e stakeholder principali. 
Gli utilizzatori sono stati suddivisi per ruolo e esigenze specifiche. È stato intervistato un totale di 14 utenti, distribuiti come segue:

\begin{itemize}
    \item \textbf{Agenti commerciali esterni:} 2 persone
    \begin{itemize}
        \item 1 agente con esperienza limitata
        \item 1 agente con esperienza consolidata
    \end{itemize}
    \item \textbf{Commerciali interni:} 4 persone
    \begin{itemize}
        \item 1 specializzato in nuovi impianti
        \item 1 specializzato in ricambi
        \item 2 con competenze trasversali (nuovi impianti e ricambi)
    \end{itemize}
    \item \textbf{Soci aziendali con responsabilità:} 3 persone
    \begin{itemize}
        \item 1 responsabile delle vendite
        \item 1 responsabile tecnico
        \item 1 responsabile amministrativo e CEO
    \end{itemize}
    \item \textbf{Direttore di produzione:} 1 persona
    \item \textbf{Impiegati tecnici:} 4 persone (non utilizzatori diretti del software)
\end{itemize}

Nel corso delle interviste è stato necessario coinvolgere utenti aggiuntivi per chiarimenti e approfondimenti su aspetti specifici.

\subsubsection{Metodologia di conduzione}
Le interviste sono state condotte in modalità mista, in presenza e da remoto, secondo la disponibilità degli intervistati.
Per gli utilizzatori abituali è stato privilegiato un approccio a distanza tramite Teams per ottenere anche una registrazione dell'utilizzo del sistema attuale.
Per gli utenti non utilizzatori è stato preferito un approccio in presenza per facilitare la comunicazione e l'interazione.

Nella fase iniziale sono state predisposte domande strutturate definite in accordo con il project manager.
Con il progredire delle interviste è stato adottato un approccio più flessibile, portando i risultati delle interviste precedenti ai nuovi incontri per approfondire aspetti specifici e chiarire dubbi emergenti.
Questo approccio adattivo ha consentito di testare proposte e direzioni progettuali in modo dinamico e di raccogliere feedback più dettagliati e pertinenti.

\subsubsection{Durata interviste}
Le interviste hanno avuto durata variabile da 30 minuti a 2 ore, in funzione della complessità delle esigenze e del livello di dettaglio richiesto.

\subsubsection{Registrazione ed elaborazione}
Tutte le interviste sono state registrate (con consenso informato) e trascritte per facilitare l'analisi successiva.
Le trascrizioni sono state analizzate per identificare i requisiti funzionali e non funzionali, le criticità del sistema attuale e le opportunità di miglioramento.

\end{document}